\section{Delayed feedback and robustness}\label{sec:delayed}

\phantomsection\label{sec:delayed-law}As a case study, the DVF command of Eq.~\eqref{eq:dvf} is replaced by
\begin{equation}\label{eq:comb-law}
  u(t)=-g_1\dot x(t)-g_2\bigl[\dot x(t)-\dot x(t-T)\bigr]-k_d\bigl[x(t)-x(t-T)\bigr],
\end{equation}
clipped as before, and studied with $k_d=0$ (DVF\,+\,D). The delayed terms vanish on every $T$-periodic motion, so the forced orbit and the free voltage $u_0$ depend on $g_1$ alone. The law is published \citep{pyragas1992continuous,vandewouw2015ifac,vandijk2015ijrnc,li2022ijamt}. Both families were designed by the authority rule of Eq.~\eqref{eq:feas} with $g_1$ up to \SI{3000}{\volt\second\per\metre}, so that in effect the same gain is compared with and without the delayed term. With $g_2=\SI{2000}{\volt\second\per\metre}$ the linear limit rises by a median factor of 1.78 (1.44--2.10) over the 251 speeds. Simulated at 25 speeds with the \SI{9}{\volt} clip, the face-entry limit rises by a median factor of 1.72, but the step-entry limit by only 1.03 (0.98--1.27), because the entry transient, not the damping of the chatter mode, saturates the actuator (Supplementary Material, Section~\ref{S-supp:std-delayed-law}).
% src: README §13.2; src/c4_comb.py (docstring, model), src/c4_robust.py (controller); tests/test_c4_comb.py (test_comb_terms_vanish_on_the_settled_cut); paper/notes/novelty.md §1.6
% src: results/c4_limits.npz (feas_g1: max |g1 - g1_DVF| = 9.7 V s/m); results/c4_summary.json (limits.feasible: d2000 ratio_to_dvf median 1.78, range 1.44-2.10)
% src: results/c4_entry.npz (face 1.45-2.12, median 1.72; step 0.98-1.27, median 1.03)

\phantomsection\label{sec:delayed-robust}Returning to DVF, the residual authority of the nominal $g=1500$ design at the four speeds of Section~\ref{sec:bistab-basin} was sampled under errors in the model: the cutting coefficients were drawn independently and uniformly within $\pm\SI{10}{\percent}$, the natural frequency within $\pm\SI{5}{\percent}$, the damping ratio within $\pm\SI{20}{\percent}$, the actuator force constant within $\pm\SI{10}{\percent}$, the eccentric runout within 0--\SI{5}{\micro\metre}, the clip level within 8.1--\SI{9}{\volt} and the loop delay within 0--\SI{300}{\micro\second}, 4000 samples, in closed form. At the design depth, $0.9$ of the nominal $\alim$, the sampled cut is linearly unstable in 8--\SI{20}{\percent} of the samples, and its forced orbit reaches the clip in up to \SI{9}{\percent}. Even at half the limit, the forced orbit reaches the clip in 69 and \SI{50}{\percent} of the samples at 2500 and \SI{2870}{\rpm}, where the nominal residual authority is only 1.5 and \SI{3.2}{\volt} at that depth; there the runout dominates (Table~\ref{S-tab:mc}).
% src: verification/paper/mc_closed_form.json (levels.pm10/pm20.ranges; n_per_case 4000; speeds.*.fracs.*.p_unstable, p_saturated, u0_stable_median_V, dominant_u0_all, spearman_u0_all: 2500 rpm 0.5 runout -0.62, 2870 rpm 0.5 runout -0.83; nominal.u0_9V); mc_closed_form.py, mc_common.py (caches restored; nominal reproduces stored a_lim and bmax: checks.nominal)
% src: verification/paper/mc_closed_form.json (nominal.<rpm>.u0_9V; levels.pm10/pm20.speeds.<rpm>.fracs.<frac>: u0_stable_median_V, p_unstable, p_saturated)
